\section{Discussion}

The decision-relevant result is a range for the number of local government areas (LGAs) above the 50\% threshold, not a single point estimate. Severe food insecurity showed positive state-level spatial association (Moran's \(I=0.204\), permutation \(p=0.018\)), whereas the moderate-or-severe measure did not (\(I=-0.012\), \(p=0.427\)). This contrast is consistent with geographic structure concentrated in the severe tail rather than shared by the broader outcome. For agency use, we therefore report the identified range of 87 to 231 LGAs, not 115 as a fixed tally. The lower endpoint nevertheless depends on model-smoothed prevalences for observed LGAs, so it is not a direct-data count.

Under the fitted joint model, the posterior expected count was \(E[C]=114.975\) of 774 LGAs, with a 95\% posterior interval from 101 to 130 and a posterior standard deviation of 7.383. That interval describes a joint distribution of threshold crossings. The marginal exceedance probabilities alone permit \(\operatorname{Var}(C)\) from 0.0244 to 5492.63, equivalent to a standard-deviation range from 0.156 to 74.11; the fitted variance of 54.509 lies only 0.99\% of the way up this feasible range. Thus, the interval from 101 to 130 is a property of the fitted dependence structure, not something identified by the marginals. The marginals fix the expected count, but they do not say whether areas cross the line together. Area-specific decisions were also uneven: 32 LGAs had an exceedance probability above 0.95, 509 were classified below the line, and 40 lay in the 0.40--0.60 probability band.

The fitted dependence contribution is not negligible. The independent variance term was 40.258, while the pairwise term was 14.248, giving a pairwise-term ratio of 0.354. Under an exchangeable area effect, the ratio was 0.187. This comparison is consistent with roughly half of the pairwise contribution reflecting shared coefficient uncertainty and half reflecting geographic dependence. The pairwise estimate itself has Monte Carlo variation: split-half estimates were 13.37 and 15.12. The decomposition therefore helps explain why the count varies more than it would under independence, but it does not establish that the fitted spatial dependence is the only plausible joint structure.

The partial-identification result shows why the expected count should not be treated as an observed total. The survey includes at least one cluster in 654 LGAs, while 120 LGAs have none. Those unclustered areas are 15.5\% of all LGAs but contain 11.04\% of the population (\(p=0.88963\)). The model-based count over the observed 654 areas spans 87 to 111, while allowing each unclustered area to contribute either zero or one adds 120 LGAs to the width of the full 144-LGA interval. In other words, 120 of the 144-LGA width rests on an assumption about places without sampled households. The observed-only share of people living in a crisis LGA was 13.72\%, and the corresponding identified interval for the population share was 10.63\% to 24.90\%. These are different decision quantities: the area count gives every LGA equal weight, while the population share weights areas by their residents. The interval is a posterior interval for the sample-analogue identified set, conditional on the model for observed LGAs and the population denominators. Its lower endpoint is model-based rather than a direct-data bound.

The validation results determine how far the area estimates can be carried. When one LGA was withheld while its neighbours remained in the fit, the design met all three pre-declared criteria at exposure 0.82, with coverage of 0.940. The matched over-exposed design at exposure 1.67 had coverage of 0.900 but failed the spread criterion, with a standardized-error standard deviation of 1.181. The whole-region holdout failed all three criteria at exposure 5.4, with coverage of 0.775. Agreement with direct estimates also depended on the holdout: the correlation was 0.975 in sample, 0.768 under the single-area design, and 0.195 when a whole region was withheld. In the latter design, predictions remained near 0.13 to 0.28 across areas with different observed outcomes. The contrast supports interpolation among sampled neighbours, not extrapolation to a region without local observations. High in-sample agreement therefore does not justify treating the fitted surface as portable across regions.

For the 120 unclustered LGAs, the exact-exposure design is the more relevant check. At exposure 0.95, unadjusted coverage was 0.890 and the standard deviation of standardized errors was 1.178, so the design failed its coverage and spread criteria. Group-level coverage ranged from 0.878 to 0.900, and one group of 40 LGAs carried the pass-or-fail verdict. Widening an individual-area interval by \(k=1.178\) raised pooled coverage to 0.940. The out-of-fold coverages for that widening factor were 0.875, 1.000, 0.925, 0.975, and 0.950. An agency using an interval for an unclustered LGA should report the widened interval and describe its observed coverage in these tests, about 0.89 before widening and 0.94 after. It should not label that interval a nominal 95\% interval. The variation across groups also means that the pooled coverage should not be read as a stable guarantee for every unsampled area.

The area-covariate analysis is informative because it separates a lack of measurable signal from a lack of usable geographic data. Three of eight candidates could not be interpolated adequately from the 1,380 sample points: nightlights had a cross-validated interpolation correlation of 0.522, travel time 0.614, and livestock 0.535. The five covariates that could be carried to the area level were largely inert after household composition was accounted for: 19 of 20 coefficient intervals contained zero. These results do not show that economic activity, market access, or pastoral livelihoods are unrelated to food insecurity. Rather, the available measurements of those dimensions could not be recovered at this resolution, while the area variables that were recoverable added no clear signal in this model. A future analysis would need better geographic measurements before interpreting a weak area-level association as evidence against these mechanisms.

The literature comparison is narrow by design. Chacón-Montalván et al. (2025) report threshold-exceedance probabilities and hotspot counts, so neither quantity is presented here as a new use of spatial food-insecurity models. Their paper maps a continuous surface in one municipality, with eight census sectors as its finest administrative unit. Nigerian estimates already exist for single local government areas from cross-sectional surveys (Omuemu, Otasowie and Onyiriuka, 2012; Kayode and Alabi, 2025) and for the 109 senatorial districts from survey indices with hotspot classification (Jerumeh, 2024); we found in that work no interval for an area estimate and no interval for the count. What we did not find, and what this paper supplies, is a set of estimates for the national set of 774 local government areas together with an interval for the number of areas above a threshold. Related estimators for area-level proportions include bivariate spatial models for small-area proportions (Yu, Xu and Yu, 2025) and direct-assisted unit-level models for rare-event prevalence (McGovern, Wilson and Wakefield, 2026); in what we could read of this work we found no interval for the number of areas above a line. The contribution here is the sharp bound for the count, the pre-declared spatial validation protocol, and the identification statement for areas without a sampled cluster, applied to Nigerian LGAs; it is not a new estimator.

Several limits define what the results do not claim. The model provides no validated interval for the count of areas above the threshold. The whole-region holdout does not support regional extrapolation. The 120 LGAs without clusters do not have a calibrated map, and the likelihood does not use survey design weights. The area-level calibration tests do not repair these gaps: they evaluate predictions for particular holdout designs, not the joint count distribution or the unobserved prevalence in every unsampled LGA. The 120 areas cannot be scored against a known truth. The lower endpoint of the identification interval also depends on model-smoothed prevalence in observed areas, so the interval does not remove model dependence from the observed part of the count.

Two further sources of uncertainty affect the population interpretation. LGA assignment used geoBoundaries polygons rather than the survey's own administrative boundaries, and population denominators are 2020 projections based on the 2006 census. The main inferential limitation, however, is the unweighted likelihood. The unweighted national prevalence was 27.74\%, compared with the survey-weighted estimate of 29.10\%, a difference of about 1.4 percentage points. The LGA surface and its count therefore describe the unweighted sample quantity, not a design-weighted population prevalence.
